\section{Eliminar la L y Software 3.0: no se trata de cambiar de nombre, sino de cambiar la ruta de los datos}

Las secciones anteriores trataron por separado la evolución de la pila de software y la eliminación del componente Language; esta sección examina ambos temas en conjunto. Liu Xianming (刘先明) recuerda en varias ocasiones que muchos términos nuevos hacen, en esencia, algo similar: modelos más grandes, más datos, menos reglas manuales y un paso más directo de la entrada a la acción. Lo que se denomina Software 3.0, VLA, VLM o de extremo a extremo no busca crear conceptos, sino encontrar rutas de datos que admitan más scaling.

\subsection{Por qué las estructuras intermedias reaparecen una y otra vez}

Esta sección explica una paradoja de ingeniería: las estructuras intermedias suelen agregarse para mejorar la controlabilidad, pero cuando el sistema alcanza cierto grado de complejidad, terminan limitando el techo del sistema de aprendizaje.

En los sistemas de conducción autónoma, las estructuras intermedias resultan atractivas: resultados de percepción, líneas de carril, cuadros delimitadores de objetos, trayectorias candidatas, go point, meta action, language token. Hacen que el sistema sea más interpretable, más fácil de depurar (debug) y más acorde con la división tradicional del trabajo de ingeniería. Sin embargo, cada estructura intermedia que se agrega puede introducir un cuello de botella de diseño manual: los datos se comprimen en algún formato definido por personas y el modelo solo puede aprender dentro de ese formato.

\begin{knowledgebox}{Ventajas y desventajas de las estructuras intermedias}
\begin{tabularx}{\textwidth}{p{0.22\textwidth}p{0.34\textwidth}X}
\toprule
Dimensión & Ventaja & Costo \\
\midrule
Interpretabilidad & Las personas entienden la salida de cada módulo, lo que facilita localizar problemas & Puede sacrificar información continua original. \\
División del trabajo & Percepción, predicción y planificación y control pueden avanzar en equipos separados & Las interfaces entre módulos se vuelven el techo del sistema. \\
Puesta en producción segura & Es más fácil aplicar reglas de protección y respaldo humano & Cuantas más reglas, más lento es el scaling de datos. \\
Aprovechamiento de datos & Las etiquetas intermedias permiten aprendizaje supervisado & Las etiquetas manuales limitan la escala y la capacidad expresiva. \\
\bottomrule
\end{tabularx}
\end{knowledgebox}

\subsection{El verdadero objetivo de eliminar la L: que los datos hablen por sí mismos}

A continuación se retoma la decisión de Xiaopeng (小鹏): eliminar la L no busca aparentar radicalidad, sino permitir que datos de mayor escala, más originales y más continuos entren en la cadena de entrenamiento.

El objetivo de eliminar el Language Bottleneck no es volver el sistema más misterioso, sino permitir que los sensores, las trayectorias de conducción y la retroalimentación de las acciones entren al entrenamiento con menos transformaciones manuales. El lenguaje puede usarse para instrucciones y semántica de alto nivel, pero si cada acción de bajo nivel tiene que pasar por language token, discretiza un mundo continuo y obstaculiza el aprendizaje autosupervisado y el aprovechamiento de datos a gran escala.

\begin{importantbox}{Definición mínima de Software 3.0}
En estas notas, Software 3.0 designa el uso de modelos grandes, grandes volúmenes de datos, señales autosupervisadas y fábricas de modelos en la nube para llevar las tareas del mundo físico de sistemas de reglas manuales a sistemas de aprendizaje con scaling sostenible.
\end{importantbox}

\subsection{Resumen de la sección}

Lo esencial de Software 3.0 no es la terminología, sino que la ruta de los datos sea más corta. Cuantas menos capas intermedias manuales haya, más posibilidades tiene el modelo de aprender directamente de datos reales; pero cuantas menos capas intermedias haya, más se necesitan evaluación, seguridad y gobernanza de ingeniería más sólidas.
